\begin{definition}
Пусть $H(\C)$~--- гильбертово пространство над полем $\C$,
$A \in \L(H) : A^* = A$, то есть $(Ax, y) = (x, Ay)\ \forall x, y \in H$.
Такой оператор $A$ называется \textit{самосопряжённым}.
\end{definition}

\noindent\textbf{Основная идея:} $A: H \to H$, будем вместо него рассматривать квадратичную форму:
\[ K\!: H \to \C \]
\[K(x)=(Ax, x)\]

\begin{theorem}[11.1]\label{11.1}
$H(\C)$, $A$~--- самосопряжённый оператор.
\begin{enumerate}
    \item $\forall x\: K(x) := (Ax, x) \in \R$
    \item Если $\lambda$~--- собственное значение оператора $A$, 
    то $\lambda \in \R$
    \item Пусть $\lambda_1 \neq \lambda_2$~--- собственные значения, 
    $e_1, e_2$~--- собственные векторы, соответствующие им. 
    Тогда $(e_1, e_2) = 0$.
\end{enumerate}
\end{theorem}

\begin{proof}
    \begin{enumerate}
        \item В силу самосопряженности $A$: $(x, Ax)=(Ax, x)=\overline{(x, Ax)}$. Число совпадает с сопряжением, значит оно вещественное.
        
        \item $\lambda$~--- собственное значение, $e$~--- собственный вектор.
        
        С другой стороны, 
        $(Ae, e) = (\lambda e, e) = \lambda (e, e)=\lambda \|e\|^2$. Но из пункта 1
        $(Ae, e) \in \R$
    
        \item $e_1, e_2, \lambda_1, \lambda_2$. 
        \[ (\lambda_1 e_1, e_2) = (Ae_1, e_2)=(e_1, Ae_2) = (e_1, \lambda_2 e_2)\]
        В силу пункта 2 $\lambda_1, \lambda_2$~--- вещественные, значит можно вынести из скалярного произведения
        \[ (\lambda_1 - \lambda_2) (e_1, e_2) = 0\]
        По условию собственные значения не равны, значит $(e_1, e_2)=0$
    \end{enumerate}
\end{proof}

\begin{unnecessary}
    \begin{statement}
        Пусть \(H( \C)\) --- гильбертово пр-во, \(A \in \L(H)\), т. ч. \(\forall x \in H\: \left(Ax, x \right) \in \R\). Тогда \(A\) --- ССО.
    \end{statement}
    \begin{proof}
        Рассмотрим произвольные \(x, y \in H\). Очевидно, что и \((x+y)\), и \((x + iy)\) тоже будут лежать в \(H\). Обозначим комплексные числа, являющиеся результатами следующих скалярных произведений как
        \[(Ax, y) = a + ib \quad (x, Ay) = c + id \quad a, b, c, d \in \R\]
        Покажем что \(a=c\) и \(b=d\):
        \begin{itemize}
            \item[ $\text{[}b = d\text{]}\!\!\!$ ]
            \(\R \ni \left(A(x + y), (x + y) \right) = \left(Ax, x \right) + \left(Ax, y \right) + \left(Ay, x \right) + \left(Ay, y \right)\).
            Первый и последний члены являются действительными числами по условию. Тогда  \(\left(Ax, y \right) + \left(Ay, x \right) = \left(Ax, y \right) + \overline{\left(x, Ay \right)} = a + ib + c - id \in \R\). Значит, \(b=d\).
            \item[$\text{[}a = c\text{]}\!$]
            \(\R \ni \left(A(x + iy), (x + iy) \right) = \left(Ax, x \right) + \left(Ax, iy \right) + \left(Aiy, x \right) + \left(Aiy, iy \right) = \left(Ax, x \right) - i\left(Ax, y \right) + i\left(Ay, x \right) + i(-i)\left(Ay, y \right)\). Первое и последнее слагаемое являются действительными числами по условию. Посмотрим на остальные: \( -i(a + ib) + i(c - id) = -ai + b + ci + d\). Это выражение должно быть действительным числом, поэтому \(a=c\).
        \end{itemize}
    \end{proof}
\end{unnecessary}
